\section{Discussion and Limitations}
\label{sec:discussion}

\subsection{Implications}
\paragraph{For adopters.} A QR report is best read as a structured screening aid with an evidence trail, not as a verdict. The split into Q and R tells an instructor whether a weak resource needs better content or better packaging, and the per-metric evidence shows where to look. The illustration suggests that even widely used titles leave room to improve: no scored title reached the excellent band, and the Adopt decision for five titles depends on the floor rule of Equation~\ref{eq:decision} as much as on the total.

\paragraph{For authors and publishers.} Several of the weakest scores are cheap to fix. Metadata was the weakest metric in the illustration, and structured descriptive and license metadata can be added without changing a word of content. Stated learning outcomes and accessibility statements are the next most tractable. The score thus points to specific, ordered improvements, which is what makes it more useful than a star rating.

\paragraph{For repositories and policy.} A catalog that published the ten scores, with evidence and version, alongside each record would let users filter on the properties they care about. A shared instrument also gives policy bodies something to measure progress against, which a recommendation alone does not provide \cite{unesco2019}.

\paragraph{For researchers.} The program in Section~\ref{sec:validation} treats rubric design as an empirical problem, as the review of existing rubrics urges \cite{yuan2015}. The most informative findings may be negative: metrics that do not agree across reviewers, or that do not predict learning.

\subsection{Limitations}
\begin{itemize}[leftmargin=*,itemsep=3pt]
  \item \textbf{One model family, no human reviewer.} Each title was scored in two passes by AI agents of the same model family, with no human scoring the same resource. The scores can carry that model's biases, and the two passes may share errors, so the agreement between them is probably an upper bound on human agreement. With ten titles we cannot estimate agreement per metric. The scores are therefore illustrations of the instrument, not measurements of the books.
  \item \textbf{Sampling.} Metrics 1 to 5 were scored from three chapters per book (first, middle and last numbered). The first pass used a single chapter, and its totals differed from the three-chapter totals by 3.5 points on average and by 9 for one title. Three chapters is still a small sample of a book with 16 to 47 chapters, and the first-pass rule (the fifth chapter) was itself ambiguous for a book with an unnumbered introduction.
  \item \textbf{Convenience sample.} The ten books were chosen by judgment, include six from one publisher, and are all in STEM or social-science introductory subjects with English text. The sample cannot support claims about OER in general.
  \item \textbf{Limited access for the reviewer.} Agents could read pages and public APIs but did not render pages in a screen reader, test contrast, import LMS packages or check every page for a license line. Some pages (for example a publisher's accessibility statement) were not readable without a script engine. Each agent listed what it could not verify, and its confidence labels are self-reports. In the first pass one agent reported that a shared working folder was overwritten by another agent and re-fetched its files; we have not checked whether this affected its readings, and in the second pass each agent had its own folder.
  \item \textbf{Scripted census scope.} The 101-book census is the author's own corpus of MkDocs textbooks, mostly generated by one pipeline, and the checks encode that pipeline's conventions (for example, phrase detection for learning objectives). The low pass rates for some checks say more about convention-matching than about the books.
  \item \textbf{Equal weights and a total.} Equal weights are a default, not a finding. The companion component-level specification keeps its two headline numbers separate on purpose \cite{oerevaluator}, whereas the QR report adds them; whether a combined total helps or misleads adopters is untested.
  \item \textbf{Author's interest.} The author wrote the rubric, the companion book and the evaluator skill, which is a conflict of interest for any claim that the instrument works. The independent studies of Section~\ref{sec:validation} are the remedy. The name ``OpenBook'' has not been searched for conflicts with existing marks or projects.
  \item \textbf{Not a learning measure.} The score examines the resource, not its effect on learners, and it is a snapshot of one date, 31 October 2026.
\end{itemize}

\subsection{Open design questions}
The illustration points to questions the specification should settle before validation. Should a NonCommercial or share-alike clause cap the license metric at a stated number? Should scripted checks determine fixed shares of some metrics, so that agents can no longer drift on them? Should a resource that fails a floor metric be reported with a warning flag even when its total is high? Each is an empirical question for Studies 1 and 2.
